Results follow the SEME axes. The primary positive contrast is constrained versus unconstrained signed concordance. Continuous high \(r\) at one locked checkpoint is reported secondarily. Prior-only quantifies the circularity budget for binary signs.

\subsection{Controls: unconstrained, prior-only, and linear imitation}

\textbf{Matched unconstrained twin (primary contrast).} On the locked fold-0 comparison against the constrained training signature, unconstrained agreement was 3/10 on TUH (\(r=-0.309\)), 4/10 on OSF (\(r=-0.331\)), and 4/10 on P-ADIC (\(r=-0.327\)). Constrained agreement on the same cohorts is 10/10 (Section~3.2). The unconstrained CAUEEG \(\mu_{\mathrm{base}}\) probe reached mean ROC-AUC 0.539 (below constrained 0.579 and below theta/alpha 0.675). Table~\ref{tab:ablation-direction} summarizes the locked direction contrast.
\begin{table}[!t]
\centering
\caption{Locked unconstrained battery versus constrained direction agreement against the constrained training signature (locked fold-0 seed-42 checkpoints). Unconstrained $r$ values are from that locked battery, not from the $\alpha_c$-sweep fold-0 re-simulation. Prior-only signed effects also reach 10/10 on signs with lower mean effect-vector correlation ($r{=}0.636$).}
\label{tab:ablation-direction}
\setlength{\tabcolsep}{4pt}
\begin{tabular}{lrrrrr}
\toprule
Dataset & $n$ & Constr.\ agree & Constr.\ $r$ & Unconstr.\ agree & Unconstr.\ $r$ \\
\midrule
TUH & 200 & 10/10 & 0.908 & 3/10 & -0.309 \\
OSF & 92 & 10/10 & 0.851 & 4/10 & -0.331 \\
P-ADIC & 145 & 10/10 & 0.918 & 4/10 & -0.327 \\
\bottomrule
\end{tabular}
\end{table}


\textbf{Prior-only baseline (circularity budget).} A constant literature-signed 2185-D effect, without a CVAE, yielded 10/10 signs at mean \(r=0.636\) (Donepezil 0.572; Memantine 0.700). Binary signs are therefore partly recoverable from the prior alone. SEME treats this as an explicit circularity budget: the CVAE is not claimed to uniquely explain 10/10 signs.

\textbf{Linear imitation control.} A ridge map trained on packed baselines, frozen drug embeddings, and disease labels to predict locked CVAE-simulated drug-minus-baseline vectors recovered 10/10 signs on TUH/OSF/P-ADIC with \(r=0.851\), \(0.639\), and \(0.775\). At the same fold-0 comparison, the locked constrained CVAE remained higher in continuous \(r\) (\(0.908\), \(0.851\), \(0.918\)). This is imitation of CVAE-generated targets, not independent biological validation.

\subsection{Direction concordance and continuous-\(r\) sensitivity}

Locked fold-0 seed-42 constrained correlations were \(r=0.908\) (TUH, \(n=200\)), \(0.851\) (OSF, \(n=92\)), and \(0.918\) (P-ADIC, \(n=145\)), each with secondary signs 10/10 (Table~\ref{tab:direction-agreement}; Fig.~\ref{fig:direction}).
% Table 2 -- Drug-response direction agreement
% Caption placement: ABOVE the tabular (Elsevier convention).
% Requires: \usepackage{booktabs}
\begin{table}[!t]
\centering
\caption{Checkpoint-specific external signature concordance at the locked fold-0
seed-42 checkpoint. Band-level direction agreement is scored between each
external cohort and the training-cohort constrained signature on five feature
blocks $\times$ two drugs (Donepezil, Memantine; maximum 10/10).
Correlation $r$ is the Pearson correlation of mean drug-minus-baseline effect
vectors (\texttt{effect\_magnitude\_correlation}). These values are not
architecture-wide performance estimates; fold and training-seed sensitivity
are reported in the Results. CAUEEG is not included in the locked
10/10 direction claim.}
\label{tab:direction-agreement}
\setlength{\tabcolsep}{4pt}
\begin{tabular}{@{}llllp{4.6cm}@{}}
\toprule
Dataset & N & Agree & $r$ & Assessment role \\
\midrule
  TUH & 200 & 10/10 & 0.908 & 1st external assessment \\
  OSF & 92 & 10/10 & 0.851 & 2nd assessment (independent) \\
  P-ADIC & 145 & 10/10 & 0.918 & 3rd assessment (acq.\ change) \\
\bottomrule
\end{tabular}
\end{table}

\begin{figure}[!t]
\centering
\includegraphics[width=\columnwidth]{figures/fig2_direction_agreement.png}
\caption{Locked fold-0 seed-42 checkpoint only: external signature concordance scored against the training constrained signature. Bars show full-sample effect-vector correlation $r$ for TUH ($r{=}0.908$, $n{=}200$), OSF ($r{=}0.851$, $n{=}92$), and P-ADIC ($r{=}0.918$, $n{=}145$). These bars are not representative of five-fold or multi-seed performance. Secondary signed agreement was 10/10 on each cohort. These values are checkpoint-specific signature concordance, not independent pharmacodynamic validation.}
\label{fig:direction}
\end{figure}
Subject-bootstrap 95\% percentile intervals were 0.874-0.900, 0.790-0.837, and 0.871-0.904. Five-fold means were \(0.385\pm0.308\), \(0.371\pm0.294\), and \(0.372\pm0.310\); signs were 10/10 on three folds, 9/10 on fold 4, and 8/10 on fold 2. Holding folds fixed and varying train seeds \(\{42,7,21,123,2024\}\), non-42 continuous \(r\) typically ranged 0.08-0.43 while signs remained mostly 9/10-10/10 (Table~\ref{tab:multiseed-direction}).
% Table 6 -- Fixed-fold multi-seed training sensitivity (READ-ONLY from aggregate_results)
\begin{table}[!t]
\centering
\caption{Fixed-fold training-seed sensitivity of external direction metrics.
Fold assignment and fold-0 subject membership were frozen at \texttt{split\_seed}${}={}$42;
only \texttt{train\_seed} varied. Seed 42 is the locked historical checkpoint
(\texttt{checkpoint\_constrained.pt}); seeds 7, 21, 123, and 2024 are independent
training-seed perturbations under the same fold split. Values are Pearson $r$
(mean Donepezil/Memantine effect-vector correlation) and sign agreement.
Source: \texttt{models/validation/multiseed\_robustness/aggregate\_results.csv}.}
\label{tab:multiseed-direction}
\setlength{\tabcolsep}{3.5pt}
\begin{tabular}{@{}rlccc@{}}
\toprule
Seed & Role & TUH ($r$; agree) & OSF ($r$; agree) & P-ADIC ($r$; agree) \\
\midrule
42 & Locked fold-0 & 0.908; 10/10 & 0.851; 10/10 & 0.918; 10/10 \\
7 & Alternate train seed & 0.408; 10/10 & 0.408; 10/10 & 0.428; 10/10 \\
21 & Alternate train seed & 0.154; 10/10 & 0.160; 10/10 & 0.106; 8/10 \\
123 & Alternate train seed & 0.147; 9/10 & 0.230; 9/10 & 0.080; 9/10 \\
2024 & Alternate train seed & 0.163; 10/10 & 0.182; 10/10 & 0.146; 10/10 \\
\midrule
\multicolumn{2}{@{}l}{Mean $r$ (all five)} & 0.356 & 0.366 & 0.336 \\
\multicolumn{2}{@{}l}{Median $r$} & 0.163 & 0.230 & 0.146 \\
\multicolumn{2}{@{}l}{SD $r$} & 0.328 & 0.288 & 0.354 \\
\bottomrule
\end{tabular}
\end{table}

Constrained signs exceeded unconstrained signs in all 15 seed\(\times\)cohort cells without establishing continuous-\(r\) stabilization. An equal-weight fold ensemble (\(r=0.642/0.626/0.627\)) did not clearly exceed prior-only mean \(r=0.636\). Under SEME, signed concordance is the more reproducible external component; locked high \(r\) is not architecture-level.

\subsection{Magnitude escalation}

Magnitude \(\bar{s}\) was tested in sequence (Table~\ref{tab:diagnostic-results}; Fig.~\ref{fig:diagnostic}).
% Table 3 -- Diagnostic magnitude discrimination (all contrasts)
% Caption ABOVE tabular (Elsevier). Requires booktabs, adjustbox.
\begin{table*}[!t]
\centering
\footnotesize
\setlength{\tabcolsep}{3.5pt}
\caption{Twin latent drug-response magnitude discrimination across escalating
external cohorts. Status strings are taken from the validation JSON
(\texttt{FAIL}, \texttt{PASS\_WITH\_CAVEATS}). The OSF primary row uses
\texttt{honest\_verdict.status}; OSF disease-label sensitivity rows have no
status field (shown as --). The Kruskal--Wallis row reports $H$ instead of
Cohen's $d$ and Kruskal $p$ in the $p$ column.}
\label{tab:diagnostic-results}
\begin{adjustbox}{max width=\textwidth}
\begin{tabular}{@{}llp{3.6cm}lll@{}}
\toprule
Dataset & Contrast & N (A / B) & $p$ & $d$ & Status \\
\midrule
  \multicolumn{6}{@{}l@{}}{\textbf{TUH}} \\
  \midrule
  TUH & Abnormal vs normal & 131 / 69 & $0.272$ & -0.13 & FAIL \\
  \multicolumn{6}{@{}l@{}}{\textbf{OSF}} \\
  \midrule
  OSF & AD vs HC & 80 / 12 & $p{<}0.001$ & 0.44 & PASS$^{\ast}$ \\
  OSF & AD vs HC (disease${=}0$) & 80 / 12 & $p{<}0.001$ & 0.64 & -- \\
  OSF & AD vs HC (disease${=}1$) & 80 / 12 & $p{<}0.001$ & 0.34 & -- \\
  \multicolumn{6}{@{}l@{}}{\textbf{P-ADIC}} \\
  \midrule
  P-ADIC & AD vs HC & 49 / 96 & $0.968$ & -0.07 & FAIL \\
  P-ADIC & AD vs HC (disease${=}0$) & 49 / 96 & $0.527$ & -0.16 & FAIL \\
  P-ADIC & AD vs HC (disease${=}1$) & 49 / 96 & $0.702$ & -0.14 & FAIL \\
  \multicolumn{6}{@{}l@{}}{\textbf{CAUEEG}} \\
  \midrule
  CAUEEG & Dementia vs Normal & 291 / 436 & $0.516$ & -0.00 & FAIL \\
  CAUEEG & Dementia vs MCI & 291 / 395 & $0.877$ & 0.07 & FAIL \\
  CAUEEG & MCI vs Normal & 395 / 436 & $0.363$ & -0.07 & FAIL \\
  CAUEEG & AD-tagged Dem.\ vs Normal & 214 / 436 & $0.624$ & -0.01 & FAIL \\
  CAUEEG & Kruskal--Wallis (3 groups) & 436 / 395 / 291 & $0.637$ & $H{=}0.90$ & FAIL \\
  CAUEEG & Dem.\ vs Normal (disease${=}0$) & 291 / 436 & $0.411$ & -0.04 & FAIL \\
  CAUEEG & Dem.\ vs Normal (disease${=}1$) & 291 / 436 & $0.776$ & 0.00 & FAIL \\
\bottomrule
\end{tabular}
\end{adjustbox}
\vspace{0.3em}\\
{\scriptsize $^{\ast}$\texttt{PASS\_WITH\_CAVEATS} (HC $n{=}12$; exploratory).}
\end{table*}

\begin{figure*}[!t]
\centering
\includegraphics[width=0.92\textwidth]{figures/fig3_diagnostic_null.png}
\caption{Forest plot of Cohen's $d$ for primary twin magnitude contrasts on TUH, OSF, P-ADIC, and CAUEEG. Approximate 95\% intervals use the Hedges--Olkin standard-error formula from reported $d$ and group sizes. Exact $n$, $p$, and $d$ values are listed in Table~\ref{tab:diagnostic-results}.}
\label{fig:diagnostic}
\end{figure*}
 TUH abnormal versus normal was null (p = 0.272, d = -0.13). OSF AD versus HC was an exploratory positive with HC \(n=12\) caveats (p = 0.000278, d = 0.44). P-ADIC AD versus HC was null (p = 0.968, d = -0.07). CAUEEG Dementia versus Normal was null (p = 0.516, d = 0.00, \(n=727\)), despite clear theta/alpha slowing on the same features (p < 0.001, d = 0.47) and good scale/PCA overlap (Fig.~\ref{fig:pca}).
\begin{figure}[!t]
\centering
\includegraphics[width=\columnwidth]{figures/fig4_pca.png}
\caption{PCA of 2185-D features: training cohort versus CAUEEG. Dashed circle marks the training 95\% radius; overlap fraction within that radius is 0.611.}
\label{fig:pca}
\end{figure}
Secondary CAUEEG contrasts and disease-label ablations remained null. Magnitude does not provide reliable external diagnosis in this escalation.

\subsection{Encoding bake-off (classical features versus latent)}

On CAUEEG Dementia versus Normal (\(n=727\)), SEME's Encoding axis compares classical feature references to frozen-twin \(\mu_{\mathrm{base}}\) readouts (Fig.~\ref{fig:encoding}; Fig.~\ref{fig:bakeoff}; Table~\ref{tab:encoding-results}). A logistic probe on \(\mu_{\mathrm{base}}\) (disease bit fixed to 0) reached mean ROC-AUC 0.579 (permutation p = 0.010). Theta/alpha reached 0.675 (OOF 0.673) and a packed 2185-D logistic reference reached 0.699 (OOF 0.699; feature-space reference, not a CVAE score). Nested heads peaked at OOF AUC 0.593 (permutation p = 0.005); MLP and boosting did not close the gap.
\begin{figure*}[!t]
\centering
\includegraphics[width=0.95\textwidth]{figures/fig5_encoding_analysis.png}
\caption{SEME Encoding bake-off on CAUEEG Dementia versus Normal ($n{=}727$). Amber: classical theta/alpha and related feature references. Teal/navy: frozen-twin $\mu_{\mathrm{base}}$ readouts. Gray band: latent-probe permutation null 5th--95th percentiles (0.449--0.559). Best nested head OOF AUC $= 0.593$. A packed 2185-D logistic feature-space reference reaches mean-fold AUC $= 0.699$ (not a CVAE score).}
\label{fig:encoding}
\end{figure*}
% Table 4 -- Encoding analysis (pairs with Fig. 5)
% Caption ABOVE tabular. Requires booktabs, adjustbox.
\begin{table}[!t]
\centering
\footnotesize
\setlength{\tabcolsep}{3pt}
\caption{CAUEEG Dementia vs Normal encoding analysis ($n = 727$).
Out-of-fold (OOF) ROC-AUC and balanced accuracy are reported.
Permutation $p$-values are included where computed (original latent probe:
200 label shuffles of 5-fold CV; best head: 200 shuffles of nested CV).
The 2185-D logistic row is a raw feature-space diagnostic reference, not a CVAE result,
not a theoretical upper bound, and not a supervised deep 2185-D classifier. The final column is OOF AUC minus the $\theta/\alpha$
pooled OOF AUC ($0.673$; mean-fold $\theta/\alpha$ AUC $= 0.675$ is used in the prose
and Fig.~\ref{fig:encoding}). Complements Fig.~\ref{fig:encoding}.}
\label{tab:encoding-results}
\begin{adjustbox}{max width=\columnwidth}
\begin{tabular}{@{}lccccc@{}}
\toprule
Model & Input & AUC & Bal.\ acc. & Perm.\ $p$ & vs.\ $\theta/\alpha$ \\
\midrule
  2185-D logistic (feature ref.) & packed 2185-D & 0.699 & 0.642 & -- & +0.026 \\
  $\theta/\alpha$ probe (ref.) & $\theta/\alpha$ & 0.673 & 0.625 & -- & +0.000 \\
  Untuned logistic & $\mu_{\mathrm{base}}$ & 0.576 & 0.568 & $0.010$ & -0.097 \\
  Tuned L2 logistic & $\mu_{\mathrm{base}}$ & 0.593 & 0.573 & $0.005$ & -0.080 \\
  MLP & $\mu_{\mathrm{base}}$ & 0.543 & 0.543 & -- & -0.130 \\
  HistGradientBoosting & $\mu_{\mathrm{base}}$ & 0.556 & 0.548 & -- & -0.117 \\
  Tuned LogReg ($\mu$+logvar) & $\mu$+logvar & 0.587 & 0.554 & -- & -0.085 \\
\bottomrule
\end{tabular}
\end{adjustbox}
\end{table}

\begin{figure}[!t]
\centering
\includegraphics[width=\columnwidth]{figures/fig_bakeoff_encoding.png}
\caption{Compact SEME Encoding bake-off panel: classical feature references versus latent/head readouts on CAUEEG Dementia versus Normal ($n{=}727$). Values match locked probe/head artifacts; DeLong annotation uses the non-locked Phase~C OOF comparison when available.}
\label{fig:bakeoff}
\end{figure}
Paired DeLong tests confirmed latent below theta/alpha (z = -3.63, p = 0.000285) and below 2185-D (z = -4.64, p $<$ 0.001), and best head below both references (z = -3.32 and -4.09; p $<$ 0.001); theta/alpha versus 2185-D was not significant (z = -1.11, p = 0.268). Unconstrained latent AUC (0.539) did not exceed constrained (0.579). The bake-off shows encoding-level attenuation under tested readouts, not an information-bottleneck proof.

\subsection{Constraint-weight sweep}

An exploratory fold-0 \(\alpha_c\) sweep (Fig.~\ref{fig:alpha}) showed signs at 10/10 by \(\alpha_c=0.05\), TUH \(r\) peaking near \(\alpha_c=0.1\) (\(r=0.897\)), and non-monotone latent AUC. It does not support a simple direction-versus-diagnosis trade-off and is not a nested \(\alpha_c\) selection.

\begin{figure}[!t]
\centering
\includegraphics[width=\columnwidth]{figures/fig6_constraint_strength.png}
\caption{Exploratory fold-0 constraint-weight ($\alpha_c$) sweep: signed direction agreement and continuous effect-vector correlation, with CAUEEG latent-probe AUC. Not a nested selection of $\alpha_c$.}
\label{fig:alpha}
\end{figure}

\subsection{ChemBERTa versus padded one-hot}

A matched fold-0 constrained retrain replaced ChemBERTa with 384-D padded one-hot drug IDs without touching the locked checkpoint. One-hot reached 10/10 signs on TUH/OSF/P-ADIC (\(r=0.176/0.194/0.174\)); rematched ChemBERTa reached 9/10 (\(r=0.104/0.110/0.076\)). Both continuous-\(r\) values sit in the low multi-seed regime. CAUEEG latent mean AUC was 0.540 (ChemBERTa rematch) versus 0.559 (one-hot). In this two-drug setting, chemical embeddings are not required for constraint-consistent signs; the locked ChemBERTa twin remains the primary historical checkpoint.
